\documentclass[12pt,letterpaper]{article}

% --- Paquetes de configuración regional y tipográfica ---
\usepackage[utf8]{inputenc}
\usepackage[spanish,es-tabla]{babel}
\usepackage[margin=2.5cm]{geometry}
\usepackage{amsmath,amsfonts,amssymb}
\usepackage{booktabs}
\usepackage{array}
\usepackage{xcolor}
\usepackage{listings}
\usepackage{caption}
\usepackage{fancyhdr}
\usepackage{hyperref}
\usepackage{float}
\usepackage{microtype}
\usepackage{tcolorbox}

% --- Colores institucionales y temáticos ---
\definecolor{codegreen}{rgb}{0.13,0.55,0.13}
\definecolor{codegray}{rgb}{0.45,0.45,0.45}
\definecolor{uctblue}{rgb}{0.08,0.22,0.45}
\definecolor{boxbg}{rgb}{0.96,0.97,0.99}

% --- Cajas destacadas ---
\tcbset{
    estudiobox/.style={
        colback=boxbg,
        colframe=uctblue,
        coltitle=white,
        fonttitle=\bfseries,
        boxrule=1pt,
        arc=3mm,
        left=4mm,
        right=4mm,
        top=3mm,
        bottom=3mm
    }
}

% --- Configuración de encabezados y pies de página ---
\pagestyle{fancy}
\fancyhf{}
\fancyhead[L]{\small\color{gray}Sistemas Inteligentes (INFO1157) -- Proyecto \#2}
\fancyhead[R]{\small\color{gray}Arquitectura MQTT y Bots de Telegram}
\fancyfoot[C]{\thepage}
\renewcommand{\headrulewidth}{0.5pt}

% --- Hipervínculos ---
\hypersetup{
    colorlinks=true,
    linkcolor=uctblue,
    citecolor=uctblue,
    urlcolor=uctblue
}

\begin{document}

% =========================================================================
% PORTADA
% =========================================================================
\begin{titlepage}
    \centering
    \vspace*{1.5cm}
    {\Large\textbf{UNIVERSIDAD CATÓLICA DE TEMUCO}}\\[0.3cm]
    {\large FACULTAD DE INGENIERÍA}\\[0.2cm]
    {\large DEPARTAMENTO DE INGENIERÍA CIVIL INFORMÁTICA}\\[2.5cm]
    
    {\LARGE\bfseries INFORME TÉCNICO}\\[0.4cm]
    {\Large\bfseries PROYECTO \#2: SISTEMA DISTRIBUIDO DE TELEMETRÍA AMBIENTAL}\\[0.5cm]
    {\large Arquitectura MQTT con Mosquitto, Procesamiento Multivariable en Ventana Móvil y Despacho Automatizado a Bots de Telegram}\\[2.5cm]
    
    \textbf{Asignatura:} Sistemas Inteligentes (INFO1157)\\[0.2cm]
    \textbf{Docente:} Dr. Alberto Caro S.\\[0.2cm]
    \textbf{Unidad:} Unidad 2 -- Sistemas de Sensores e Internet de las Cosas (IoT)\\[2.0cm]
    
    \begin{flushright}
        \textbf{Estudiante:} Ingeniería Civil Informática\\
        \textbf{Repositorio:} \url{https://github.com/Diegood16/Sistemas2}\\
        \textbf{Fecha:} Noviembre de 2026\\
        \textbf{Temuco, Chile}
    \end{flushright}
    \vfill
\end{titlepage}

\newpage
\tableofcontents
\newpage

% =========================================================================
% 1. RESUMEN Y OBJETIVOS
% =========================================================================
\section{Resumen Ejecutivo y Objetivos}

El monitoreo de la calidad del aire y variables meteorológicas es un pilar fundamental en las ciudades inteligentes (\textit{Smart Cities}). Este proyecto implementa un ecosistema distribuido de telemetría ambiental basado en el protocolo industrial \textbf{MQTT} (\textit{Message Queuing Telemetry Transport}), integrando captura de datos periódica, intermediación por broker centralizado, análisis en ventana móvil de 10 muestras y generación automatizada de reportes gráficos despachados a teléfonos móviles mediante la API de Telegram.

\subsection{Objetivos Principales}
\begin{enumerate}
    \item \textbf{Capa de Transporte:} Desplegar el broker MQTT Mosquitto en \texttt{127.0.0.1:1883} para coordinar el intercambio desacoplado de mensajes.
    \item \textbf{Capa de Emisión:} Implementar un nodo publicador periódico que transmita mediciones de material particulado (\texttt{DATA/MP}) y variables meteorológicas (\texttt{DATA/MA}) a razón de 5 segundos.
    \item \textbf{Capa de Suscripción y Buffering:} Construir cuatro clientes suscriptores que capturen el flujo de datos y mantengan un buffer temporal de 10 muestras para cada magnitud física.
    \item \textbf{Capa de Renderizado y Visualización:} Replicar la especificación con cinco subplots apilados en \texttt{Matplotlib}, aplicando una arquitectura orientada a objetos inmune a colisiones de concurrencia.
    \item \textbf{Capa de Notificación y Ciberseguridad:} Conectar el sistema a cuatro bots independientes de Telegram utilizando variables de entorno protegidas (\texttt{.env}) para resguardar las credenciales en repositorios públicos.
\end{enumerate}

% =========================================================================
% 2. FUNDAMENTOS CONCEPTUALES: ¿CÓMO FUNCIONA MQTT?
% =========================================================================
\section{Fundamentos Conceptuales: Protocolo MQTT y Arquitectura IoT}

\subsection{El Patrón Publicador/Suscriptor (Pub/Sub)}
En la web clásica (HTTP), la comunicación se basa en el modelo cliente-servidor directo: un cliente pide un recurso y el servidor responde. Si múltiples terminales quisieran recibir datos de un sensor ambiental mediante HTTP, el sensor debería responder peticiones individuales recurrentes, agotando su capacidad de cómputo y ancho de banda.

El patrón \textbf{Pub/Sub} de MQTT resuelve esto mediante \textbf{desacoplamiento total}:
\begin{itemize}
    \item \textbf{El Publicador (Emisor):} No requiere conocer la identidad ni la cantidad de receptores existentes. Publica la información bajo una etiqueta temática denominada \textbf{Tópico} (por ejemplo, \texttt{DATA/MP}).
    \item \textbf{El Suscriptor (Receptor):} Registra ante el broker los tópicos de su interés. Al publicarse información en dicho canal, el cliente la recibe de forma inmediata.
    \item \textbf{El Broker (Servidor Central):} Actúa como nodo intermediario. Recibe los paquetes emitidos y los retransmite únicamente a los clientes con suscripción activa.
\end{itemize}

\subsection{El Broker Eclipse Mosquitto}
\textbf{Eclipse Mosquitto} es un servidor (\textit{broker}) de código abierto altamente optimizado que implementa las especificaciones del protocolo MQTT. 
\begin{itemize}
    \item \textbf{Eficiencia:} Implementado en C, presenta una huella de memoria mínima y alta capacidad para gestionar conexiones concurrentes.
    \item \textbf{Puerto estándar:} Opera en el puerto TCP \texttt{1883} para comunicaciones en red local sin cifrado adicional.
\end{itemize}

\begin{tcolorbox}[estudiobox, title=Principio de Desacoplamiento en MQTT]
\textbf{Funcionamiento Conceptual:} El nodo emisor despacha periódicamente lecturas hacia el broker bajo un tópico determinado (como \texttt{"DATA/MP"}). Mosquitto no manipula la carga útil; inspecciona su tabla de suscripciones activas y reenvía una copia exacta a cada suscriptor registrado. De este modo, emisores y receptores operan de manera totalmente asíncrona e independiente.
\end{tcolorbox}

\subsection{Estructura de Tópicos y Cargas Útiles (Payloads)}
En este proyecto se utilizan dos tópicos con cadenas en formato de texto plano separado por comas (CSV):

\begin{enumerate}
    \item \textbf{Tópico \texttt{DATA/MP} (Material Particulado):}
    \begin{equation*}
        \text{Formato: } \texttt{"MP01,MP25,MP10\textbackslash n"} \quad \longrightarrow \quad \text{Ejemplo: } \texttt{"7, 28, 35"}
    \end{equation*}
    Representa partículas en suspensión en microgramos por metro cúbico ($\mu\text{g/m}^3$): ultrafinas ($1.0\,\mu\text{m}$), finas ($2.5\,\mu\text{m}$) y torácicas ($10\,\mu\text{m}$).
    
    \item \textbf{Tópico \texttt{DATA/MA} (Medio Ambiente):}
    \begin{equation*}
        \text{Formato: } \texttt{"Temperatura,Humedad\textbackslash n"} \quad \longrightarrow \quad \text{Ejemplo: } \texttt{"18, 72"}
    \end{equation*}
    Representa la temperatura ambiental en grados Celsius ($^\circ\text{C}$) y la humedad relativa en porcentaje ($\%$).
\end{enumerate}

% =========================================================================
% 3. FLUJO DE EJECUCIÓN PASO A PASO
% =========================================================================
\section{Flujo de Ejecución del Sistema: Del Sensor a la Notificación}

La transmisión y procesamiento de datos a través de la infraestructura se estructura en cinco etapas coordinadas:

\begin{enumerate}
    \item \textbf{Generación y Emisión (Publicador):} El script \texttt{publicador.py} se activa cada 5 segundos. Selecciona de forma balanceada si despachar lecturas de partículas o de clima, y las publica hacia el broker en \texttt{127.0.0.1:1883}.
    \item \textbf{Enrutamiento Central (Broker Mosquitto):} Mosquitto recibe el paquete y lo replica de inmediato a los cuatro clientes suscriptores que mantienen conexiones TCP abiertas y activas.
    \item \textbf{Acumulación en Ventana Móvil (Suscriptores):} Cada cliente suscriptor recibe la cadena de texto en su función \texttt{on\_message()}, la descompone con \texttt{split(',')} y guarda los números en listas de memoria.
    \item \textbf{Disparo del Gráfico (10 Muestras):} Cuando el cliente detecta que acumuló al menos 10 lecturas de partículas y 10 de clima, toma las primeras 10 muestras y activa la función \texttt{generar\_grafico()}.
    \item \textbf{Envío por Telegram (Despacho a Teléfono):} El gráfico generado se almacena como imagen PNG de alta fidelidad. Luego, utilizando el token de su bot asignado y el \textit{Chat ID} del usuario, invoca el método \texttt{sendPhoto()} de la API de Telegram, notificando al usuario en su dispositivo móvil.
\end{enumerate}

% =========================================================================
% 4. MAPEO DE MÓDULOS DEL CÓDIGO
% =========================================================================
\section{Estructura del Software y Organización Modular}

El proyecto fue diseñado de manera modular para asegurar la separación clara de responsabilidades:

\begin{table}[H]
\centering
\caption{Mapeo de archivos y responsabilidades del Proyecto \#2.}
\label{tab:modulos}
\begin{tabular}{llp{7.5cm}}
\toprule
\textbf{Archivo} & \textbf{Ubicación} & \textbf{Responsabilidad Principal} \\
\midrule
\texttt{.env} & \texttt{Proyecto2/} & Credenciales privadas (Tokens y Chat ID) ignoradas por Git. \\
\texttt{.env.example} & \texttt{Proyecto2/} & Plantilla pública de configuración sin tokens reales. \\
\texttt{config.py} & \texttt{Proyecto2/src/} & Lee el archivo \texttt{.env}, define puertos MQTT y credenciales de los bots. \\
\texttt{publicador.py} & \texttt{Proyecto2/src/} & Emite lecturas periódicas cada 5 segundos. \\
\texttt{subscriber.py} & \texttt{Proyecto2/src/} & Clase base: gestión de red MQTT, ventana de 10 muestras, gráficos y Telegram. \\
\texttt{cliente1.py} al \texttt{4.py} & \texttt{Proyecto2/src/} & Scripts ejecutables independientes para lanzar cada cliente por separado. \\
\texttt{simulador\_completo.py} & \texttt{Proyecto2/src/} & Ejecutable todo-en-uno que levanta hilos paralelos para toda la solución. \\
\bottomrule
\end{tabular}
\end{table}

\subsection{Detalle del Funcionamiento de los Cuatro Clientes}
Cada cliente suscriptor representa una estación o bot temático especializado:
\begin{itemize}
    \item \textbf{Cliente 1:} Asociado al \textbf{Bot 1} (Foco: Material Particulado PM 1.0).
    \item \textbf{Cliente 2:} Asociado al \textbf{Bot 2} (Foco: Material Particulado PM 2.5).
    \item \textbf{Cliente 3:} Asociado al \textbf{Bot 3} (Foco: Material Particulado PM 10).
    \item \textbf{Cliente 4:} Asociado al \textbf{Bot 4} (Foco: Medio Ambiente, Temperatura y Humedad).
\end{itemize}
Cada uno recibe ambos tópicos para contextualizar la gráfica completa de 5 subplots, despachando su reporte hacia su bot respectivo con su propio token individual.

% =========================================================================
% 5. INGENIERÍA Y RESOLUCIÓN DE PROBLEMAS TÉCNICOS
% =========================================================================
\section{Decisiones de Ingeniería y Desafíos Resueltos}

Durante la implementación surgieron desafíos técnicos que fueron resueltos aplicando principios sólidos de ingeniería de software.

\subsection{Gestión de Concurrencia en Matplotlib}
Al ejecutar los cuatro clientes suscriptores en paralelo, se identificó un riesgo de colisión al renderizar las figuras de manera simultánea.
\begin{itemize}
    \item \textbf{Causa Técnica:} La interfaz estándar \texttt{matplotlib.pyplot} opera sobre una máquina de estados global compartida que no está diseñada para acceso concurrente multi-hilo. Llamadas concurrentes a \texttt{subplots()} o \texttt{close()} entre diferentes hilos pueden reiniciar el contexto de dibujo durante la exportación a disco.
    \item \textbf{Solución Implementada:} Se migró la generación gráfica a la \textbf{API Orientada a Objetos de Matplotlib}:
    \begin{itemize}
        \item Se prescindió del uso de \texttt{pyplot} en el módulo de generación.
        \item Cada hilo instancia de manera independiente un objeto \texttt{Figure(figsize=(7.5, 7))} con su propio motor de dibujo \texttt{FigureCanvasAgg(figura)}.
        \item Se incorporó un bloqueo de exclusión mutua (\texttt{threading.Lock}) y validación de tamaño de archivo previa al despacho por red.
        \item \textbf{Resultado:} Generación estable de reportes de alta calidad sin interferencias entre procesos concurrentes.
    \end{itemize}
\end{itemize}

\subsection{Ciberseguridad y Manejo de Credenciales con \texttt{.env}}
Para resguardar los tokens de autenticación de los bots de Telegram se implementaron directrices estándar de seguridad:
\begin{itemize}
    \item Los tokens se alojan exclusivamente en el archivo local \texttt{.env}.
    \item El archivo \texttt{.gitignore} restringe la indexación de archivos de entorno (\texttt{*.env}) y directorios de caché de Python (\texttt{\_\_pycache\_\_}).
    \item Se incluye \texttt{.env.example} como plantilla de referencia con los nombres de variables requeridos sin datos confidenciales.
    \item El cargador de configuración valida la existencia de credenciales antes del envío por Telegram, emitiendo mensajes de diagnóstico informativos en la consola.
\end{itemize}

% =========================================================================
% 6. CONCLUSIONES
% =========================================================================
\section{Conclusiones}

El desarrollo del Proyecto \#2 demostró con éxito la aplicabilidad del protocolo MQTT para la supervisión y telemetría de magnitudes físicas en tiempo real. La arquitectura construida logra un equilibrio óptimo entre:

\begin{itemize}
    \item \textbf{Escalabilidad:} La separación de roles a través del Broker Mosquitto permite incorporar decenas de nuevos sensores o clientes sin alterar el código existente.
    \item \textbf{Robustez Concurrente:} La migración a la API Orientada a Objetos de Matplotlib resolvió de raíz las condiciones de carrera entre hilos simultáneos.
    \item \textbf{Ciberseguridad Práctica:} El aislamiento de credenciales en \texttt{.env} asegura que el proyecto pueda ser distribuido y auditado en cualquier computadora sin riesgos de seguridad ni filtración de tokens privados.
\end{itemize}

\end{document}
